% !TEX root = ../main.tex
\chapter{앞선 연구 살펴보기와 이론의 틀}
\label{ch:literature}

\ref{ch:introduction}장에서 세운 연구 문제는 다음과 같아요. 송금(transfer) 그래프(점과 화살표로 이은 그림)로 매기는 평판 점수는, 토큰을 대신 써도 된다는 허락에 아무 역할도 주지 않아요. 이 허락은 믿음을 분명하게 드러내는 행동이고, 모든 ERC-20 토큰(토큰 공통 규칙을 따르는 토큰)이 이것을 기록해요. 또 이런 점수의 그래프는 돈을 낼 때마다 커져요. 목표 O1--O5와 연구 질문 RQ1--RQ5는 이 문제에서 나와요. 이 문제를 두고 엔도스랭크(EndorseRank)와 적응형 가중 페이지랭크(AWP)를 견주려고, 우리는 각 방법의 그래프에 페이지랭크(PageRank)를 돌려요. 그리고 그렇게 나온 순위를 순위 상관(두 줄 세우기가 얼마나 비슷한지 나타내는 값)으로 점검해요. 이론의 틀은 하이퍼링크(웹 페이지끼리 잇는 링크)--인용 퍼뜨림(전파), 분야 일치도(점수가 재려는 분야와 얼마나 잘 맞는지), 계산 복잡도(계산하는 데 드는 수고가 얼마나 빨리 늘어나는지)를 합친 것이에요. 남은 연구의 빈틈(아직 밝혀지지 않은 부분)은 \ref{ch:methodology}장과 \ref{ch:results}장에서 다뤄요.

\dissertationSubheading[sec:blockchain-defi-foundations]{블록체인과 디파이의 기초}

누구나 볼 수 있는 공개 블록체인은 거래를 해시(내용을 짧은 지문 같은 값으로 바꾼 것)로 묶은 블록에 기록해요. 그리고 각 블록은 그 뒤에 덧붙는 블록들이 지켜 줘요\cend{5}. 참여하는 컴퓨터(노드) 가운데 절반이 넘는 수가 한 블록을 받아들이고 나면, 그 블록을 고쳐 쓰는 데 큰 비용이 들어요. 이 비용 덕분에, 거래하는 사람들은 저마다 알려진 법적 신분이 없어도 장부 위에서 거래를 마무리할 수 있어요(\emph{믿음을 최소로 줄인} 결제).

비트코인(Bitcoin)은 이런 장부 위에 사람과 사람이 직접 주고받는(P2P) 전자 현금 네트워크를 만들었어요\cend{5}. 이더리움(Ethereum)도 뒤에 덧붙이기만 할 수 있는 장부를 그대로 써요. 하지만 계정이 잔액과 함께 프로그램도 가질 수 있게 했어요. 그러면 거래 하나가 그 프로그램의 함수(정해진 일을 하는 명령 묶음) 하나를 불러서 계약의 상태를 바꿀 수 있어요\cend{7}. 이 함수들은 이더리움 가상 머신(EVM, 이더리움 안에서 프로그램을 돌리는 장치)에서 돌아가요\cend{36}. 이렇게 장부에 저장된 프로그램을 스마트 계약(블록체인 위에서 저절로 실행되는 프로그램)이라고 불러요. 정직한 노드라면 어느 것이든, 같은 함수 호출을 다시 돌려 보면 똑같은 결과를 얻어요\cend{6}. 그래서 담보(돈을 빌릴 때 맡겨 두는 재산) 비율, 청산(손실이 너무 커져 거래가 강제로 끝나는 일) 조건, 수수료 표를 믿을 만한 제삼자 대신 계약이 지키게 할 수 있어요\cend{1}.

이런 계약들이 모여 디파이(DeFi, 탈중앙 금융)를 이루어요. 이 계약들은 함께 거래소, 대출 서비스, 파생상품(다른 자산의 값에 따라 값이 정해지는 상품) 시장, 자산 관리자가 하는 일을 해요. Werner 등\cend{8}은 디파이 서비스(프로토콜)의 종류를 정의했고, Sch\"{a}r\cend{1}는 서비스들이 서로의 위에 어떻게 층층이 쌓이는지 설명했어요. 한 서비스가 다른 서비스와 맞물려 움직일 수 있어서, 이미 있는 금융 상품들을 엮어 새 금융 상품을 만들 수 있어요\cend{37}. 그런데 바로 이 성질 때문에 한 곳의 고장이 여러 서비스로 번질 수도 있어요. 실제로 성공한 공격에는 플래시 론(한 거래 안에서 빌리고 바로 갚는 대출)을 악용한 공격\cend{38}과 가격 오라클(블록체인 밖의 가격을 알려 주는 장치) 조작\cend{39}이 있었어요. 이런 블록체인의 거래 기록은 누구나 볼 수 있고 컴퓨터가 읽을 수 있어요. 하지만 그 기록은 행동을 담을 뿐, 법적 신분은 담지 않아요\cend{2}.

\dissertationSubsubheading[sub:pseudonymity-credit]{가명성과 신용 문제}

보통의 신용 점수 매기기는 신용 기록을 법적 주체(법으로 인정받는 사람이나 회사)에 묶어요. 그런데 공개 블록체인에서는 분석의 기본 단위가 주소(지갑)예요. 주소 하나는 법적 주체일 수도 있고, 신탁(남의 재산을 맡아 관리하는 단체)일 수도 있고, 알고리즘(정해진 계산 순서)으로 움직이는 프로그램이나 스마트 계약일 수도 있어요. 한 주체가 주소를 여러 개 가질 수 있고\cend{11}, 통계로 주소를 묶어 보아도(군집화) 누가 주인인지는 일부만 알아낼 수 있어요\cend{40}. 그래서 온체인(블록체인 위의) 대출에서는 빌리는 사람이 빌린 돈보다 값이 더 나가는 담보를 맡겨야 해요. 담보의 값이 정해 둔 선 아래로 떨어지면, 청산하는 사람이 빚의 일부를 대신 갚고 그만큼의 담보를 싸게 받아 갈 수 있어요\cend{41}. 이 방식은 안전하지만, 돈을 묶어 둬요. 온체인 신용 점수는 블록체인 위의 행동만 보고, 어떤 참여자가 돈을 갚을 가능성이 얼마나 되는지를 수로 나타내려고 해요. 최근 제안들은 전망 이론(사람이 이익과 손해를 어떻게 느끼는지 설명하는 이론)\cend{13}, 대출 자료에 쓴 기계 학습\cend{14}, 계정 연결망을 따라 신용 위험이 퍼지는 모형\cend{12}을 바탕으로 해요.

\dissertationSubheading[sec:ethereum-l2]{이더리움, 레이어 2 롤업, 아비트럼 원}

대부분의 디파이 앱은 이더리움 메인넷(이더리움의 본래 블록체인)에 올라가 있어요. 이더리움 메인넷은 지금까지 디파이의 새로운 시도와 유동성(쉽게 사고팔 수 있도록 시장에 모인 돈)이 가장 많이 모인 곳이었어요. 하지만 그곳을 쓰는 사람이 늘면서 가스(거래 수수료) 비용이 오르고, 한 번에 처리할 수 있는 양이 모자라게 되었어요\cend{42}. 레이어 2(L2) 방식들은 이더리움의 안전 보장은 지키면서 거래를 더 싸게 만들려고 해요. 어떤 것은 거래를 오프체인(블록체인 밖의) 공간에서 실행하고, 어떤 것은 거래를 묶어서 실행 결과를 짧게 줄인 요약을 이더리움 메인넷에 올려요\cend{43}. Thibault 등\cend{44}은 L2 규모 키우기 방법들을 종류별로 나눈 분류 체계로 정리했어요. 아비트럼(Arbitrum) 같은 낙관적 L2(일단 맞다고 보고, 문제가 있을 때만 따지는 L2)는 제출된 거래 묶음을, 이의 제기 기간 안에 반대 증거가 나오지 않으면 올바른 것으로 봐요\cend{45}.

아비트럼 원(Arbitrum One)은 이런 낙관적 L2 가운데 하나이고, EVM에 맞춰 만든 스마트 계약을 실행할 수 있어요. 아비트럼 원은 디파이 앱들이 즐겨 쓰는 곳이 되었고, 그 가운데에는 무기한 선물 거래(끝나는 날 없이, 값이 오를지 내릴지에 거는 거래)를 하는 거래소인 GMX도 있어요\cend{46}. 우리가 아비트럼 원을 쓰는 까닭은, 공개된 빅쿼리(BigQuery, 큰 자료를 찾아보는 서비스) 자료에 아비트럼 원의 ERC-20 \texttt{Approval} 이벤트와 \texttt{Transfer} 이벤트(계약이 남기는 기록)가 모두 들어 있기 때문이에요. 또 수수료가 충분히 낮아서, 여섯 달 동안의 기록에도 활동이 빽빽하게 들어 있기 때문이에요. 더 검증하지 않고서는, 우리 결과가 다른 L2나 이더리움 메인넷에도 똑같이 들어맞는다고 주장하지 않아요. 하지만 우리가 제안한 틀(허락 그래프와 송금 그래프, 구성개념(점수가 재려고 하는 생각) 점검, 순위 상관)은 같은 로그를 남기는 블록체인이라면 어디에나 쓸 수 있다고 믿어요.

\dissertationSubheading[sec:erc20-allowance-lifecycle]{ERC-20 토큰과 허락의 한살이}

ERC-20 표준은 대체 가능한 토큰(하나하나가 똑같은 값어치여서 서로 바꿔도 되는 토큰)을 위한 공통 인터페이스(프로그램끼리 서로 약속한 함수 모음)를 정해요. 이 인터페이스는 \texttt{transfer}, \texttt{approve}, \texttt{transferFrom}, \texttt{allowance}와 여기에 딸린 이벤트들로 이루어져요\cend{18}. 주인(토큰 주인)이 \texttt{approve(spender, amount)}를 부르면, 토큰 계약은 스펜더(spender, 허락을 받아 대신 쓰는 쪽)가 주인의 잔액에서 \texttt{amount}까지 옮겨도 된다고 기록해요. 그러면 스펜더는 \texttt{transferFrom}을 불러 토큰 송금을 실제로 할 수 있어요. EIP-2612 서명 허락(\texttt{permit})을 쓰면, 오프체인 서명 하나로 허락(allowance, 토큰을 얼마까지 대신 써도 된다는 허락)을 바꿀 수 있고, 이 일은 거래 한 번으로 끝나요\cend{19}. 쓰는 사람은 더 편해지지만, 허락이 뜻하는 바는 그대로예요.

그림~\ref{fig:allowance-sequence}에 ERC-20 허락 하나의 한살이(처음부터 끝까지 거치는 과정)를 차례로 나타냈어요. 허락은 써도 된다는 권한을 주는 일이고, 실제 송금은 일어나지 않아요. 토큰 주인은 라우터(거래를 대신 이어 주는 계약), 금고(vault), 거래 계약에 허락을 줄 수 있어요. 이것들은 주인이 정해 둔 한도 안에서만 움직이리라고 믿는 계약이에요. 허락 거두기(0을 허락하기)와 바꾸기(0이 아닌 새 허락 정하기)는 둘 다 지금 효력이 있는 허락을 바꿔요. 엔도스랭크는 이 점을 반영해서, 주인--스펜더 쌍마다 가장 최근의 허락 값을 화살표(엣지)의 무게(가중치)로 써요. 그래서 새 허락이 앞의 허락을 통째로 바꿔요. 옛 값이 시간이 지나며 조금씩 줄어드는 것이 아니에요.

\begin{figure}[htbp]
\centering
\begin{tikzpicture}[
  font=\small,
  actor/.style={draw, rounded corners, minimum width=2.4cm, minimum height=0.75cm, align=center},
  msg/.style={-{Stealth[length=2.5mm]}, thick},
  note/.style={font=\scriptsize\itshape, text=black!70, align=left}
]
% Actors
\node[actor] (owner) at (0,0) {주인};
\node[actor] (token) at (5.5,0) {ERC-20 토큰};
\node[actor] (spender) at (11,0) {스펜더};

% Lifelines (vertical dashed lines)
\draw[dashed, gray] (owner.south) -- ++(0,-6.6);
\draw[dashed, gray] (token.south) -- ++(0,-6.6);
\draw[dashed, gray] (spender.south) -- ++(0,-6.6);

% Message 1: approve (owner -> token), horizontal span 5.5cm
\draw[msg] (0,-1.2) -- node[above,font=\scriptsize] {\texttt{approve(spender,$a$)}} (5.5,-1.2);
\node[note] at (2.75,-1.7) {\texttt{Approval} 기록; 허락 저장};

% Alternative to message 1: EIP-2612 permit, submitted by the spender (or a relayer) with the owner's signature
\draw[msg, densely dotted] (11,-2.6) -- node[above,font=\scriptsize] {\texttt{permit(owner,spender,$a$,sig)}} (5.5,-2.6);
\node[note] at (8.25,-3.1) {\texttt{approve} 대신: 주인이 오프체인에서 서명;\\ \texttt{Approval} 기록; 허락 저장};

% Message 2: transferFrom (spender -> token), horizontal span 5.5cm
\draw[msg] (11,-4.2) -- node[above,font=\scriptsize] {\texttt{transferFrom(owner,to,$v\leq a$)}} (5.5,-4.2);
\node[note] at (8.25,-4.7) {토큰 옮기기; 허락 줄이기};

% Message 3: revoke (owner -> token), dashed
\draw[msg, dashed] (0,-5.8) -- node[above,font=\scriptsize] {\texttt{approve(spender,$0$)} (허락 거두기)} (5.5,-5.8);
\end{tikzpicture}
\caption[{ERC-20 허락의 한살이예요. 허락하기(approve) 또는 서명 허락(permit), 해도 되고 안 해도 되는 \texttt{transferFrom}, 허락 거두기로 이어져요. 엔도스랭크 화살표는 주인--스펜더 쌍의 기간 안 마지막 허락을 토큰마다 더한 값이고, 그 허락은 나이와 금액으로 무게를 매겨요.}]{ERC-20 허락의 한살이예요. \texttt{approve}나 서명으로 대신하는 \texttt{permit}(점선), 해도 되고 안 해도 되는 \texttt{transferFrom}, 그리고 허락 거두기로 이어져요. 엔도스랭크 화살표 하나는 토큰마다 항을 하나씩 두고 이를 더한 값이에요. 각 항은 관찰 기간 안에서 그 주인--스펜더 쌍의 가장 최근 \texttt{Approval} 이벤트에서 나와요. 이 항은 AWP가 송금에 무게를 매기듯이, 그 허락의 나이(지난 시간)와 금액으로 무게를 매겨요. 나중 허락이 앞선 허락을 대신해요.}
\label{fig:allowance-sequence}
\end{figure}

우리가 가진 자료에서 허락은 송금에 비해 듬성듬성해요. 허락 한 번이 여러 번의 송금을 감당할 수 있고, 많은 지갑은 바깥 계정에 허락을 한 번도 해 준 적이 없어요. 이 희소성(화살표가 적음)은 계산에서 좋은 점이에요. 또 뜻의 차이도 보여 줘요. 엔도스랭크는 써도 된다는 권한을 수로 나타내지만, 송금 그래프는 돈이 흘러간 기록을 담기 때문이에요.

\dissertationSubheading[sec:graph-reputation]{그래프로 매기는 온체인 평판}

페이지랭크와 그 변형들은, 방향이 있는 화살표를 따라 점(노드)으로 흘러 들어오는 영향력을 보고 점에 순위를 매겨요\cend{15}. 최근 두 연구가 이 방법을 블록체인 거래 그래프에 썼는데, 두 연구는 서로 구별해야 해요. Do와 Do\cend{4}는 사회적 신용(다른 사람들과의 관계에서 생기는 믿음)을 재려고 이더리움 거래 그래프에 페이지랭크와 HodgeRank를 돌렸어요. 이 연구의 주된 기여(연구가 새로 보탠 것)는 신용 점수를 위해 계정 사이가 얼마나 이어져 있는지(연결성)를 정의한 것이에요. Do, Do, Nguyen\cend{3}은 적응형 가중 페이지랭크(AWP)를 제안했어요. AWP는 이 논문이 엔도스랭크와 나란히 평가하는 송금 그래프 방법이에요(\ref{sec:awp-in-depth}절). Lin 등\cend{12}은 위험을 연결망 전체로 퍼뜨리지 않으면, 거래로 이어진 모습만으로는 위험과 관련된 구조를 놓칠 수 있다고 주장해요. Wu 등\cend{24}도 피싱(속여서 돈이나 정보를 빼앗는 사기) 찾기에서 비슷한 결론에 이르렀어요. 거기서는 거래 그래프의 연결망 임베딩(그래프 속 위치를 숫자 목록으로 바꾼 것)이 계정 하나하나의 특징보다 더 잘 해냈어요. 블록체인 기반 평판 시스템에 대한 조사\cend{25}와 스마트 계약 평판 설계에 대한 조사\cend{26}에 따르면, 실제로 쓰이는 시스템 대부분은 아직도 사람이 직접 매긴 점수나 송금량을 믿음의 신호로 써요. 암호화폐 거래 자료에서 지식을 찾는 연구에 대한 조사\cend{47}는 차수(이어진 화살표 수)에 바탕을 둔 중심성과 금액에 바탕을 둔 중심성(점이 그래프에서 얼마나 중요한 자리에 있는지 나타내는 값)을 지갑을 설명하는 주된 값으로 꼽아요. Packin과 Lev-Aretz\cend{2}는 사용자가 가명으로(진짜 이름 대신 주소로만) 활동하더라도, 이런 시스템은 남이 따져서 확인할 수 있어야 한다고 강조해요.

연구 문헌 밖에서, 실제로 쓰이는 디파이와 web3(블록체인 위에서 돌아가는 인터넷 서비스) 시스템은 보통 다른 방법으로 믿음을 쌓아요. 주소는 누구나 장부에서 살펴볼 수 있는 활동으로 평판을 쌓고, 그 주소에 묶인 도구가 그 활동을 요약해요. 가장 널리 쓰이는 도구는 휴먼 패스포트(Human Passport)예요. 예전 이름은 깃코인 패스포트(Gitcoin Passport)였어요. 이 도구는 정부 신분증, 생체 정보(얼굴이나 지문 같은 몸의 정보), 다른 사람이 대신 확인해 주기, 예전의 web2와 web3 활동 같은 점검마다, 스탬프(stamp)라고 부르는 확인할 수 있는 자격 증명을 모아요. 스탬프마다 위조하기가 얼마나 어려운지를 나타내는 무게가 있고, 무게를 매겨 더한 값이 기준값에 이르면 그 주소는 통과해요\cend{64}. 패스포트는 아비트럼을 포함한 여러 네트워크의 EVM 주소를, 온체인 활동을 보고 사람인지 시빌(한 사람이 만든 가짜 주소)인지로 나누기도 해요\cend{64}. 패스포트의 목적은 시빌 공격(가짜 주소를 많이 만들어 속이는 공격)을 막는 것이에요. 그래서 온체인 증거와 오프체인 증거를 합치고, 등급을 매기는 순서가 아니라 들여보낼지 말지 정하는 문처럼 일해요. 사람 증명(주소 뒤에 진짜 한 사람이 있음을 보이는 것) 방식들을 살펴본 한 검토 연구는, 이 방식들에서 시빌 공격 막기를 늘리면 자기 주권(자기 정보를 스스로 다스리기)과 사생활 보호가 줄어든다고 말해요\cend{65}. 최근의 한 보안 검토는 값싼 스탬프를 한꺼번에 많이 얻는 자격 증명 파밍을 패스포트에 남은 가장 큰 약점으로 꼽아요\cend{66}. 이런 도구들은 한 주소를 별개의 한 사람이 다루고 있는지에 답해요. 엔도스랭크는 다른 질문에 답해요. 어떤 주소가 다른 주소들로부터 자기 토큰을 대신 써도 된다는 허락을 받는지예요. 그리고 엔도스랭크는 온체인 기록만으로 그 주소들에 순서를 매겨요(표~\ref{tab:method-comparison-lit}).

그래프 방법이 매력적인 까닭은 (i) 공개된 정보만 쓰고, (ii) 순위 매기기에 알맞은 연속 점수(끊김 없이 이어지는 값)를 내고, (iii) 화살표 수로 계산 복잡도를 분석할 수 있기 때문이에요. 하지만 그래프 방법은 시빌 공격\cend{11}과 자전 거래(자기끼리 사고팔아 거래가 많은 척하기)에 약해요. 또 그래프에서 중심에 있는 정도와 실제로 빚을 못 갚을 위험(채무 불이행 위험)이 서로 맞지 않는 문제도 있어요. 페이지랭크 점수는 금융 규제에서 말하는 신용 점수가 아니에요. 여기서의 목표는 허락으로 만든 순위와 송금으로 만든 순위를, 같은 대리 지표(대신 재는 값) 묶음에 대어 꼼꼼하게 비교하는 것뿐이에요.

\dissertationSubheading[sec:pagerank-formalism]{페이지랭크의 수식}

점이 $|V|$개인 방향이 있는 그래프 $G=(V,E)$를 생각해 봐요. $i$에서 $j$로 가는 화살표의 무게는 $w_{ij}>0$으로 써요. 이 무게는 $i$가 $j$에게 주는 영향력을 나타내요. 그 영향력은 앞에서 소개한 보증(믿고 맡김)의 뜻이어도 되고, 송금의 뜻이어도 돼요. 이에 맞는 무게 인접 행렬(어느 점이 어느 점과 얼마의 무게로 이어졌는지 적은 표)을 $W$라고 하고, 점 $i$의 나가는 무게의 합을 $o_i=\sum_j w_{ij}$라고 해요. 그러면 각 줄의 합이 1인 이동 행렬(전이 행렬) $P$는 다음과 같아요.
\begin{equation}
P_{ij} =
\begin{cases}
w_{ij}/o_i & \text{만약 } o_i>0,\\
1/|V| & \text{만약 } o_i=0
\end{cases}
\label{eq:transition}
\end{equation}
(매달린 점(나가는 화살표가 없는 점)은 자기 몫을 모든 점에 고르게 나눠 줘요). 감쇠 계수(화살표를 따라갈 확률) $d\in(0,1)$에 대해, 구글 행렬(Google matrix)은 다음과 같아요.
\begin{equation}
G_{\mathrm{PR}} = d\, P^{\top} + \frac{1-d}{|V|}\,\mathbf{1}\mathbf{1}^{\top},
\label{eq:google-matrix}
\end{equation}
그리고 페이지랭크 벡터(수를 한 줄로 늘어놓은 것) $\mathbf{r}$은 $\mathbf{r}\geq 0$과 $\sum_i r_i=1$을 지키면서 $\mathbf{r}=G_{\mathrm{PR}}\mathbf{r}$을 만족해요\cend{15,23}. $P$에서 매달린 점의 줄을 가운데 항으로 따로 떼어 쓰고 원소(벡터 속 수 하나하나)마다 풀어 쓰면, 이 식은 다음과 같아요.
\begin{equation}
r_j = d\sum_{i:(i\to j)\in E} r_i \frac{w_{ij}}{o_i} + \frac{d}{|V|}\sum_{i:\,o_i=0} r_i + \frac{1-d}{|V|}.
\label{eq:pagerank-node}
\end{equation}

\dissertationSubsubheading[sub:random-surfer]{무작위로 돌아다니는 사람으로 풀이하기}

식~\eqref{eq:pagerank-node}의 뜻은 확률로 풀어 읽을 수 있어요. 걸음마다, 나가는 화살표가 있는 점에 선 무작위로 돌아다니는 사람은 확률 $d$로 그 화살표 가운데 하나를 따라가요. 어느 화살표를 고를지는 그 무게에 비례해요. 아니면 확률 $1-d$로, 모든 점 가운데 고르게 뽑은 한 점으로 순간 이동(화살표를 무시하고 아무 점으로나 건너뛰기)해요(그림~\ref{fig:random-surfer}). 나가는 화살표가 없는 점에서는 언제나 고르게 뽑은 점으로 뛰어가요. 이것이 식~\eqref{eq:transition}의 매달린 점 규칙이에요. 그러면 정상 분포(오래 걸었을 때 각 점에 머무는 비율)는 오랫동안 걸었을 때 점마다 얼마나 자주 들르는지로 점에 순위를 매겨요. 감쇠는 순위 웅덩이(밖으로 나가는 화살표가 없는 점 무리)가 몫을 모두 빨아들이지 못하게 막아요. $d<1$이고 순간 이동이 고르게 일어나면, $G_{\mathrm{PR}}$의 모든 칸이 양수(0보다 큰 수)예요. 그래서 $G_{\mathrm{PR}}$ 행렬은 원시 행렬(몇 번 거듭 곱하면 모든 칸이 양수가 되는 행렬)이고, 그 페론 벡터(Perron vector, 여기서는 페이지랭크 벡터)는 하나뿐이에요\cend{23}. 그 증명은 부록~\ref{sec:pr-existence}에 있어요.

\begin{figure}[htbp]
\centering
\begin{tikzpicture}[
  font=\small,
  vertex/.style={draw, circle, minimum size=0.95cm},
  arrow/.style={-{Stealth[length=2.5mm]}, thick},
  teleport/.style={arrow, dashed}
]
\node[vertex] (a) at (0,0) {$A$};
\node[vertex] (b) at (3.2,1.5) {$B$};
\node[vertex] (c) at (3.2,-1.5) {$C$};
\node[vertex] (d) at (6.4,0) {$D$};
% Graph transitions: A->B, A->C, B->D, C->D (same topology as the worked example; D is dangling)
\draw[arrow] (a) -- node[above left, font=\footnotesize] {$2$} (b);
\draw[arrow] (a) -- node[below left, font=\footnotesize] {$1$} (c);
\draw[arrow] (b) -- node[below left=-1pt, font=\footnotesize] {$3$} (d);
\draw[arrow] (c) -- node[above left=-1pt, font=\footnotesize] {$1$} (d);
% Uniform teleportation, drawn from D only: one dashed arrow to every node, including D itself.
% D->A runs through the gap between B and C, and D->B / D->C bend outward, so no dashed arrow crosses a solid edge.
\draw[teleport] (d) -- (a);
\draw[teleport] (d) to[bend right=32] (b);
\draw[teleport] (d) to[bend left=32] (c);
\draw[teleport] (d) to[loop right, min distance=12mm] (d);
\end{tikzpicture}
\caption[{\ref{sec:pagerank-worked-example}절의 점 네 개짜리 그래프로, 무작위로 돌아다니는 사람의 눈으로 페이지랭크를 쉽게 풀어 본 그림이에요. 실선 화살표는 그래프를 따라 옮겨 가는 길이고, 모두 합쳐 확률 $d$로 따라가요. 점선 화살표는 고른 순간 이동이에요. 지금 있는 점을 포함해 모든 점으로 확률 $(1-d)/n$씩 가요. 점선 부채꼴은 $D$에서만 그렸지만, 사실은 모든 점에서 나가요.}]{\ref{sec:pagerank-worked-example}절의 점 네 개짜리 그래프로, 무작위로 돌아다니는 사람의 눈으로 페이지랭크를 쉽게 풀어 본 그림이에요. 화살표에 붙은 수는 허락의 무게예요. 실선 화살표는 그래프를 따라 옮겨 가는 길이에요. 돌아다니는 사람은 모두 합쳐 확률 $d$로 그 가운데 하나를 따라가고, 어느 것을 고를지는 무게에 비례해요. 점선 화살표는 고른 순간 이동을 나타내요. 순간 이동은 확률 $(1-d)/|V|$로 모든 점(지금 있는 점도 포함)에 닿아요. 점선 부채꼴은 $D$에서만 그렸지만, 사실은 모든 점에 하나씩 있어요. $D$에는 나가는 화살표가 없어요. 그래서 식~\eqref{eq:transition}의 매달린 점 규칙에 따라, $D$에 있는 사람은 언제나 고르게 뽑은 점으로 뛰어가요. 각 점으로 갈 확률은 $1/|V|=0.25$예요. 감쇠가 아니라 바로 이 규칙이, $D$가 걷기 전체를 빨아들이지 못하게 막아요.}
\label{fig:random-surfer}
\end{figure}

\dissertationSubsubheading[sub:power-iteration]{거듭제곱 반복과 계산 복잡도}

실제로는 보통 거듭제곱 반복(같은 계산을 되풀이하기)으로 $\mathbf{r}$을 구해요. 고른 벡터 $\mathbf{r}^{(0)}=\mathbf{1}/|V|$에서 시작해서 다음 계산을
\begin{equation}
\mathbf{r}^{(t+1)} = G_{\mathrm{PR}}\mathbf{r}^{(t)},
\label{eq:power-iteration}
\end{equation}
$\|\mathbf{r}^{(t+1)}-\mathbf{r}^{(t)}\|_1 < \varepsilon$이 될 때까지, 또는 반복(한 바퀴) 횟수 상한 $I_{\max}$에 이를 때까지 되풀이해요. 듬성듬성한 그래프에서는 반복 한 번에 $O(|E|)$번(화살표 수에 비례하는 만큼)의 셈이 들어요\cend{48}. 그러므로 화살표를 더 적게 만드는 방법은 반복 한 번에 드는 비용을 줄이고, 때로는 필요한 메모리도 줄여요. 송금 기록 전체 대신 마지막 허락(가장 최근 허락)의 스냅숏(어느 한 순간의 모습)을 쓰는 방법이 그 예예요. 가설 H3은 이 계산 복잡도 사실에서 나왔어요.

이 연구는 기본으로 감쇠를 흔히 쓰는 값인 $d=0.85$로 정해요\cend{15}. 그리고 이 연구에서만 정한 풀이 프로그램(solver) 설정 두 가지를 써요. 수렴(값이 더는 거의 바뀌지 않게 됨)의 허용 오차 $\varepsilon=10^{-8}$과, 최대 반복 횟수 $I_{\max}=300$이에요.\footnote{허용 오차와 반복 횟수 상한은 실행 시간을 다시 재도 같게 나오도록(재현성) 고른 값이에요. 고전적인 페이지랭크 정의의 한 부분으로 내놓은 값은 아니에요.} 강건성 점검(조건을 바꿔도 결과가 버티는지 보기) 절에서는 감쇠를 $d\in\{0.75,0.85,0.95\}$로 바꿔 가며 살펴봐요.

\dissertationSubheading[sec:reputation-systems]{평판 시스템이 이어져 온 흐름}

페이지랭크는 스펙트럼 방법(행렬의 고유벡터로 점수를 구하는 방법)과 퍼뜨림에 바탕을 둔 방법으로 평판을 매기는, 더 넓은 방법 가족에 속해요. Kleinberg의 HITS 알고리즘은 페이지마다 점수 두 개를 줘요. 하나는 허브(좋은 페이지를 많이 가리키는 페이지)로서의 점수이고, 다른 하나는 권위(좋은 페이지들이 많이 가리키는 페이지)로서의 점수예요\cend{49}. EigenTrust는 이 생각을 P2P 네트워크로 넓히고, 이웃끼리 주는 믿음을 정규화(합이 1이 되게 맞추기)해요. 그래서 나쁜 뜻을 가진 점이 제 점수를 마음대로 올릴 수 없어요\cend{16}. TrustRank는 웹 스팸(검색 순위를 속이려는 가짜 페이지)을 누르려고, 믿을 수 있는 점들로 이루어진 씨앗 집합에서 출발해 그 믿음을 바깥으로 퍼뜨려요\cend{17}. 이 모든 방법에서 평판은 관계에서 생겨요. TrustRank처럼 꼼꼼히 골라 둔 씨앗 집합에서만 믿음이 흘러나오면, 새로 온 점은 이미 평판을 가진 점들의 보증 없이는 점수를 올릴 수 없어요. 이 연구는 엔도스랭크를 AWP와 맞추려고 고른 순간 이동을 택했는데, 이렇게 하면 그 보호가 더 약해져요. 얼마나 약해지는지는 \ref{sec:sybil-cost-model}절에서 따져 봐요.

온체인에서는 여전히 시빌 공격이 가장 큰 걱정거리예요\cend{11}. 새 지갑은 거의 돈을 들이지 않고 만들 수 있어요. 하지만 서로 다른 거래 상대들에게서 허락이나 송금을 받는 지갑은 만드는 데 비용이 더 들어요. 이 연구에서 쓰는 시빌 보정 안정성 대리 지표(들어오는 거래 상대 비율, 활동 기간, 활동한 달 수)는 이런 거래의 다양함을 어느 정도 담아요. 하지만 나쁜 뜻을 가진 사람이 이 지표들을 조작하는 것을 막지는 못해요. \ref{ch:discussion}장에서는 허락 그래프에서 시빌을 만드는 데 드는 비용 모형을 내놓아요.

디파이 특징을 쓰는 기계 학습 신용 모형\cend{14}은, 라벨(정답으로 붙여 둔 값)이 달린 채무 불이행을 예측하도록(미리 맞히도록) 학습해요. 그러려고 해석할 수 있음(왜 그런 점수가 나왔는지 알 수 있음)을 포기해요. 엔도스랭크와 AWP는 해석할 수 있고 확인하기 쉬운 퍼뜨림 규칙을 지켜요\cend{2}.

\dissertationSubheading[sec:allowance-endorsement]{보증으로 보는 ERC-20 허락}

사용자는 ERC-20의 \texttt{approve} 함수와 EIP-2612의 \texttt{permit}로, 스펜더가 자기 대신 얼마까지 옮겨도 되는지를 분명하게 밝혀요\cend{18,19}. 엔도스랭크는 토큰마다 한 주인이 한 스펜더에게 준 마지막 허락을 보증 하나로 봐요. 그리고 AWP가 송금에 무게를 매기는 방식으로 이 허락에 무게를 매기고, 그렇게 만든 방향이 있는 그래프에 페이지랭크를 돌려요. 새 허락이 옛 허락을 대신하므로, 쌍과 토큰마다 이벤트가 하나씩만 들어가요.

좀 더 자세히 말하면, 주인--스펜더 쌍마다 토큰별로 기간 안의 마지막 허락이, 시간 무게(오래될수록 가벼워지는 무게)에 상한이 있는 금액 무게(금액이 아무리 커도 1을 넘지 않는 무게)를 곱한 값만큼을 보태요. AWP가 송금에 무게를 매기는 것과 같아요. 마지막 값이 0인 쌍(허락 취소)은 빼고, 남은 항들을 토큰에 걸쳐 더해요(\ref{sub:endorserank-graph}절의 식~\eqref{eq:approve-weight}). 그렇게 만든 그래프 $G_{\text{approve}}$는 보통, 같은 지갑과 같은 기간으로 만든 송금 그래프보다 훨씬 더 듬성듬성해요. 이 그래프에 페이지랭크를 돌리면 엔도스랭크 점수 $\mathbf{s}_{\mathrm{ER}}$이 나와요.

허락을 보증으로 보는 것은 일부러 한 선택이고, 꼭 들어맞지는 않는 선택이에요. 사용자가 라우터를 허락하는 것은 편해서이지, 꼭 그 라우터가 믿을 만하다고 남들 앞에서 말하려는 것은 아니에요. 또 계약 주소와 외부 소유 계정(EOA, 개인 키로 움직이는 계정)이 함께 섞여 있어요. 그래도 허락은 받은 송금보다 찬성의 뜻을 더 분명하게 보여 줘요. 받은 송금은 시장 조성(늘 사고팔 물량을 내놓는 일), 자전 거래\cend{29}, 에어드롭 파밍(공짜 토큰을 노린 가짜 활동)\cend{32}, 또는 보증할 뜻이 전혀 없는 경로 이어 주기에서 생길 수도 있기 때문이에요.

\dissertationSubheading[sub:gf-pr]{이 연구가 주장하지 않는 것}

연구 계획서에서는 두 방법을 GMX~V2 거래 성공에 비추어 검증하려고 했어요. 그리고 결과에 바탕을 둔 기준인 GF-PR로 그 검증의 윗한계를 정하려고 했어요. GF-PR은 실제로 낸 이익과 손실로 만든 별 모양 그래프에 페이지랭크를 돌린 것이에요. 그런데 GF-PR은 순환 논리(답을 이미 질문 속에 넣고 있음)예요. GF-PR은 거래 성공 대리 지표를 정하는 바로 그 이익·손실 기록으로 계산해요. 그래서 GF-PR이 실현 이익 대리 지표와 거의 완벽하게 일치하는 것은, 한 변수와 그 변수를 단조 변환(순서를 그대로 지키며 바꾸기)한 값 사이의 상관일 뿐이에요. 그리고 GF-PR과 함께, 그것에 기대던 연구 질문과 주장도 거두어들였어요. 같은 기간의 거래 대리 지표는 보관해 둔 지표 가족(비슷한 지표 묶음)으로만 남아 있어요(표~\ref{tab:alignment-summary}). 거래 결과는 기간 밖(점수를 매긴 기간 뒤)에서 다시 나와요. 매칭 코호트(모든 점수가 똑같이 점수를 매기는 지갑 묶음)의 트레이더(거래하는 사람)를 대상으로 한 탐색적(미리 정하지 않고 살펴본) 실행에서 나오고(\ref{sub:trader-results}절), 미리 등록한 6월부터 8월까지의 기간에서 청산 없이 닫은 비율(강제로 끝나지 않고 스스로 거래를 닫은 비율)로 나와요(\ref{sub:fresh-holdout}절). 이 논문의 어떤 점수도 거래 이익, 손실 피하기, 청산을 예측하는 값으로 내놓지 않아요.

\dissertationSubheading[sec:rank-correlation-theory]{순위 상관 이론}

평판 점수와 검증용 대리 지표는 둘 다 연속 변수(끊김 없이 이어지는 값)이고, 어느 쪽에도 둘로 나눌 자연스러운 경계가 없어요. 이 연구는 AWP의 검증 설계\cend{3}와 고전 심리 측정학(사람의 마음과 능력을 수로 재는 학문)\cend{22}을 따라, 두 가지 순위 일치를 재요. 하나는 단조 일치(한쪽이 커질 때 다른 쪽도 커지는지)이고, 다른 하나는 쌍 일치(둘씩 짝지었을 때 순서가 같은지)예요. 부트스트랩(다시 뽑기) 신뢰구간(믿을 만한 범위)은 보고하는 모든 계수의 표본 불확실성(표본을 뽑을 때마다 생기는 흔들림)을 수로 나타내요\cend{34}.

스피어먼의 순위 상관\cend{20}은 순위들끼리 구한 피어슨 상관(두 값이 함께 움직이는 정도)이에요. 동점(값이 같은 것)이 없는 $n$개 항목의 짝지은 순위 $(R_i,S_i)$에서는, 이 값이 다음처럼 간단해져요.
\begin{equation}
\rho = 1 - \frac{6\sum_{i=1}^n (R_i-S_i)^2}{n(n^2-1)}.
\label{eq:spearman}
\end{equation}
동점이 있으면 이 동점 없는 식은 맞지 않아요. 그때 스피어먼의 로($\rho$)는 평균 순위(가운데 순위)의 피어슨 상관이에요. 평균 순위에서는 동점인 항목들이 자기들이 차지한 자리의 평균을 함께 나눠 가져요. 이 논문에서 보고하는 $\rho$는 모두 이렇게 계산했어요. 두 식 모두에서 $\rho$는 단조 연관성을 재요.

켄달의 타우($\tau$, 두 줄 세우기가 얼마나 같은 순서인지 나타내는 수)\cend{21}는 순서가 맞는 쌍과 어긋나는 쌍의 수를 바탕으로 해요. $i<j$인 쌍 $(i,j)$는 모두 $n_0=n(n-1)/2$개예요. 그 가운데 $n_c$개는 순서가 맞는 쌍이고(두 순위가 같은 쪽으로 줄 세움), $n_d$개는 어긋나는 쌍이에요(두 순위가 반대로 줄 세움). 어느 한 순위에서라도 동점인 쌍은 둘 다 아니에요. 첫째 순위에서 $k$번째 동점 무리의 크기를 $t_k$라고 할 때, $n_1=\sum_k t_k(t_k-1)/2$이라고 해요. $n_2$도 둘째 순위에 대해 똑같이 정해요. 그러면 $\tau$-b 변형은 다음과 같아요.
\begin{equation}
\tau_b = \frac{n_c - n_d}{\sqrt{(n_0-n_1)(n_0-n_2)}},
\label{eq:kendall}
\end{equation}
동점 보정을 빼면 이 식은 $\tau_a=(n_c-n_d)/n_0$가 돼요. 이 논문에서 보고하는 $\tau$는 모두 $\tau_b$이고, 손대지 않은 원래 점수와 대리 지표로 계산했어요. 이 차이는 동점이 아주 많은 점수에서 중요해요. 매칭 코호트의 엔도스랭크가 바로 그런 점수예요(\ref{sec:rank-divergence}절). 둘씩 짝지은 순서가 얼마나 일치하는지 볼 때는, 보통 켄달의 $\tau$가 더 해석하기 쉬운 계수예요. 그래서 \ref{ch:results}장은 이 계수를 지표 가족 수준의 주된 요약으로 써요. 그 값은 같은 코호트에서 다른 방법들의 값과 견주어요. 출처를 밝히지 않은 절대 기준값은 쓰지 않아요.

\dissertationSubheading[sec:theoretical-framework]{이론의 틀}

이 연구는 엔도스랭크와 AWP를 먼저 평판 구조(H1), 구성개념 일치도(H2), 계산 효율(H3)에서 견줘요. 그다음에는 기간 밖 예측(RQ4)과 두 층(두 종류의 그래프)의 결합(RQ5)에서 견줘요. 앞부분의 틀은 세 갈래의 이론, 곧 하이퍼링크--인용 퍼뜨림, 구성개념(분야) 일치도, 계산 복잡도가 잡아 줘요. 뒷부분에는 두 가지 논리가 더 관련돼요.

\begin{itemize}
\item 하이퍼링크--인용의 뜻: 처음의 페이지랭크 알고리즘\cend{15}에서는 들어오는 하이퍼링크 하나를 인용 하나로 쳐요. Do와 Do\cend{4}는 이 논리를 이더리움 송금에 썼어요. AWP\cend{3}는 송금 화살표마다, 그 송금들이 얼마나 오래되었는지와 금액으로 무게를 매겨요. 나이에는 로지스틱 감쇠(오래된 일일수록 가볍게 세는 무게)를 쓰고, 금액에는 상한이 있는 변환을 써요(\ref{sub:awp-graph}절). 엔도스랭크는 인용에 바탕을 둔 퍼뜨림을 ERC-20의 허락하기(approve) 화살표와 서명 허락(permit) 화살표로 넓혀요. 허락은 쓸 권한을 온체인에서 분명하게 맡기는 일이고, 거두거나 바꾸기 전까지 효력이 이어져요\cend{18,19}. 반면 송금은 믿으려는 뜻 없이 그냥 거래나 경로 이어 주기를 나타낼 수도 있어요. 그래프에 바탕을 둔 위험 모형들\cend{12}은 연결망 구조가 활동을 그냥 센 수보다 더 많은 정보를 담을 수 있음을 보여 줘요. 이것은 송금량과 함께 방향이 있는 보증 그래프도 써야 한다는 근거가 돼요. 어떤 라벨에서 퍼뜨림이 그냥 센 수보다 무엇을 더 보태는지는 실제 자료로 따져 볼 문제예요. 홀드아웃(나중 기간을 떼어 두고 맞혀 보는 시험)이 이것을 바로 시험해요.

\item 분야 일치도 틀: 측정 이론은 구성개념 일치도를 써서, 어떤 지표가 따라가야 할 구성개념을 얼마나 가깝게 따라가는지 물어요\cend{22}. Messick\cend{33}은 이 생각을 넓혀서, 점수가 무엇을 뜻하는지와 그 점수를 쓰면 무엇이 뒤따르는지를 하나로 묶어 판단하게 했어요. 순위로 매기는 평판 점수에서는 스피어먼의 $\rho$와 켄달의 $\tau$가 순서의 단조 일치와 쌍 일치를 재고, AWP 검증 절차\cend{3}처럼 둘로 나누는 기준값이 필요 없어요. 그래프가 앞뒤가 맞는지(일관성)는 같은 기간의 송금 통계와 허락 통계로 점검해요. 기간 밖 점검은 시간 홀드아웃이에요. 스펜더에게는 새 허락과 새 송금자(새로 송금을 보내온 주소)를 쓰고, 미리 등록한 6월부터 8월까지의 기간에서는 트레이더에게 청산 없이 닫은 비율을 써요.

\item 계산 복잡도: 페이지랭크는 반복 한 번에 $O(|E|)$가 들어요\cend{15}. 허락 그래프는 송금 그래프보다 듬성듬성해서 화살표 수 $|E|$가 더 적어요. 그래서 반복 한 번에 드는 시간과 메모리가 적고, 실제로는 수렴하는 데 반복 횟수도 더 적게 드는 때가 많아요. 이 논리가 H3의 바탕이에요. 두 방법은 같은 풀이 프로그램을 써요. 그러므로 실행 시간의 차이는 화살표가 듬성듬성해서 생기는, 예상된 효과예요. 이것은 알고리즘의 기여로 치지 않아요.

\item 기준선(비교 기준)으로 쓰는 주변 개수(한 점에 바로 붙은 것만 센 수): 이미 거래 상대가 많은 지갑은 거래 상대를 더 많이 얻는 경향이 있어요. 그런 경우에는 원시 차수(화살표를 그냥 센 수)가 앞으로 새로 올 상대를 잘 예측해요. 그리고 퍼뜨린 점수는, 보증하는 쪽이 몇이냐보다 누구냐가 더 중요할 때에만 원시 차수를 이길 수 있어요. 그래서 RQ4는 각 페이지랭크를 동결일(점수를 고정하는 날)의 자기 원시 차수와 맞대어 봐요.

\item 층 결합: 다층 페이지랭크는 여러 관계를, 한 점 집합 위에 놓인 여러 층으로 다뤄요\cend{62,63}. C-PR(결합 페이지랭크, 두 그래프를 함께 걷는 점수)처럼 점마다 층을 섞으면, 촘촘함이 아주 다른 층들에 대해 풀리지 않은 질문이 남아요. 듬성듬성한 층에서 화살표가 없는 점은 촘촘한 층을 따라가요. 그래서 듬성듬성한 층의 신호가 살아남을 수도 있고, 사라질 수도 있어요. 따라서 RQ5는 두 가지를 물어요. 섞은 걷기가 송금 걷기보다 나은지, 그리고 두 층을 모두 이기는지예요.
\end{itemize}

그림~\ref{fig:conceptual-model}의 개념 모형은 이 논문의 점수들을 H1--H3, RQ1--RQ4와 이어 줘요. RQ5는 오른쪽의 점선 상자예요. RQ5는 H2/RQ2와 RQ4에서 입력을 받아, 두 구성개념이 하나로 합쳐지는지 물어요.

\begin{figure}[htbp]
\centering
% !TEX root = ../main.tex
\begin{tikzpicture}[
  font=\small,
  box/.style={draw, rounded corners, align=center, minimum width=10.5cm, minimum height=1.0cm},
  smallbox/.style={draw, rounded corners, align=center, minimum width=3.4cm, minimum height=1.05cm},
  refbox/.style={draw, dashed, rounded corners, align=center, minimum width=3.4cm, minimum height=0.95cm},
  arrow/.style={-{Stealth[length=2.5mm]}, thick}
]
\node[box] (center) at (0,3.2) {한 지갑 코호트에서 구한 허락 점수와 송금 점수};

\node[smallbox] (rep) at (-5,1.2) {평판\\구조};
\node[smallbox] (risk) at (0,1.2) {구성개념\\점검};
\node[smallbox] (eff) at (5,1.2) {계산\\비용};

\node[smallbox] (h1) at (-5,-1.0) {H1 / RQ1\\엔도스랭크--AWP\\갈라짐};
\node[smallbox] (h2) at (0,-1.0) {H2 / RQ2\\각 구성개념과\\맞는 정도};
\node[smallbox] (h3) at (5,-1.0) {H3 / RQ3\\점수마다\\드는 비용};

\node[refbox] (hold) at (0,-3.0) {RQ4: 시간 홀드아웃\\점수 대 원시 차수};
\node[refbox] (rq5) at (5,-3.0) {RQ5: 결합 연산자 C-PR\\화살표 두 종류, 라벨 두 개};

\draw[arrow] (center.south -| rep.north) -- (rep.north);
\draw[arrow] (center.south) -- (risk.north);
\draw[arrow] (center.south -| eff.north) -- (eff.north);

\draw[arrow] (rep.south) -- (h1.north);
\draw[arrow] (risk.south) -- (h2.north);
\draw[arrow] (eff.south) -- (h3.north);

\draw[arrow, dashed] (h2.south) -- (hold.north);
\draw[arrow, dashed] (h2.south east) -- (rq5.north west);
\draw[arrow, dashed] (hold.east) -- (rq5.west);
\end{tikzpicture}

\caption[{개념 모형이에요. 허락 점수와 송금 점수를 평판 구조(H1/RQ1), 구성개념 점검(H2/RQ2), 계산 비용(H3/RQ3), 그리고 앞으로 생길 새 허락과 새 송금자를 원시 차수와 맞대는 시간 홀드아웃(RQ4)에 이어요.}]{한 지갑 코호트의 허락 점수와 송금 점수를 평판 구조(H1/RQ1), 구성개념 점검(H2/RQ2), 계산 비용(H3/RQ3), 시간 홀드아웃에 이어 주는 개념 모형이에요. 시간 홀드아웃에서는 각 페이지랭크를, 그 페이지랭크가 매끄럽게 다듬는(평활하는) 원시 차수와 맞대어 봐요(RQ4). 오른쪽의 점선 상자는 RQ5, 곧 결합 연산자 C-PR이에요.}
\label{fig:conceptual-model}
\end{figure}

H1에는 엔도스랭크와 AWP만 들어가요. H2/RQ2와 H3/RQ3은 모든 점수에 대한 것이에요. RQ4는 스펜더 코호트에서 시간을 나누어 하는 시험이에요. 라벨은 앞으로 생길 새 주인--스펜더 허락과 새 송금자이고, 기준선은 동결일의 원시 차수예요. RQ4는 거래 결과에 순위를 매기지 않아요.

\dissertationSubheading[sub:terminology]{용어}

정식 구성개념은 \ref{ch:introduction}장에서 정의했어요. 아래 용어들은 결과를 풀이할 때 써요:

\begin{itemize}
\item 한 층 방법: 엔도스랭크(허락 그래프)와 AWP(송금 그래프)예요. 한 종류의 화살표 위만 걷는 두 페이지랭크 점수예요. \ref{sub:coupled-graph}절의 결합 걷기들과 원시 차수도 이것들과 나란히 평가해요.
\item 주된 그래프 구성개념: 한 방법이 순위를 매기도록 \emph{설계된} 그래프예요. 예를 들어 엔도스랭크에서는 $G_{\text{approve}}$예요. 점수를 매긴 뒤에 하는 점검과는 구별해요.
\item 같은 기간 점검: 점수와 같은 기간에서 뽑은, 한 점 주변의 송금 통계나 허락 통계예요\cend{22}.
\item 시간 홀드아웃: 점수를 $t_1$(2026년 2월 28일)에 고정하고, 새 주인--스펜더 허락 쌍과 새 송금자를 $t_2$(2026년 5월 31일)까지 세요. 미리 등록한 두 번째 기간에서는 2026년 5월 31일에 점수를 고정하고, 2026년 6월부터 8월까지 라벨을 세요(\ref{sub:fresh-holdout}절).
\item 매칭 지갑 코호트($n=5{,}521$), 스펜더 홀드아웃 코호트($n=1{,}335$), 확장 지갑 풀($N=99{,}080$, 더 큰 주소 모음): \ref{sub:cohorts}절의 세 주소 묶음이에요.
\end{itemize}

\dissertationSubheading[sec:domain-alignment-metrics]{분야 일치도 지표}

순위로 매기는 점수에서는 구성개념 일치도를, 점수와 점검 순위 사이의 스피어먼의 $\rho$와 켄달의 $\tau$로 재요. 이것은 AWP 검증 절차\cend{3}와 Cronbach와 Meehl\cend{22}을 따른 것이에요. 이 값은 같은 기간의 송금 점검과 허락 점검, 그리고 시간 홀드아웃의 두 라벨에 대해 재요. 실제로 어떻게 재는지에 대한 정의는 \ref{ch:methodology}장에 있어요. 이렇게 지금까지 살펴본 문헌은 세 가지 설계 선택을 뒷받침해요. 해석할 수 있는 퍼뜨림의 뼈대인 페이지랭크, 송금과 따로 떼어 낸 보증 화살표인 허락, 그리고 순위 상관에 시간 나누기를 더한 점검 설계예요.

\dissertationSubheading[sec:pagerank-worked-example]{풀어 본 예: 점 네 개짜리 보증 그래프의 페이지랭크}

지갑 네 개 $\{A,B,C,D\}$가 다음과 같은 마지막 허락 화살표로 이어져 있다고 해 봐요.
\[
A\to B\ (2),\quad A\to C\ (1),\quad B\to D\ (3),\quad C\to D\ (1),
\]
감쇠는 $d=0.85$로 해요. 나가는 무게의 합은 $o_A=3$, $o_B=3$, $o_C=1$, $o_D=0$이에요(마지막 것은 매달린 점이에요). 화살표 $B\to D$와 $C\to D$는 각각 자기 점에서 나가는 유일한 화살표예요. 그래서 각 화살표는 확률 1로 따라가게 되고, 그 무게는 점수에 아무 영향도 주지 않아요. 순간 이동은 첫 걸음을 포함해 걸음마다 모든 점에 $(1-d)/4=0.0375$를 더해 줘요. 고른 벡터 $\mathbf{r}^{(0)}=\mathbf{1}/4$에서 거듭제곱 반복을 한 걸음 하면, $(A,B,C,D)$에 대해 $\mathbf{r}^{(1)}=(0.0906,\,0.2323,\,0.1615,\,0.5156)$이 나와요. 벌써 몫이 $D$ 쪽으로 기울어요. $B$와 $C$가 둘 다 $D$를 보증하기 때문이에요. 반면 $A$는 자기 몫을 $B$와 $C$에 나눠 줘요. 반복은 정상 벡터 $\mathbf{r}=(0.1375,\,0.2154,\,0.1765,\,0.4706)$으로 수렴해요. 그래서 $D$가 맨 앞에 서고, $B$와 $C$가 그 뒤를 잇고, $A$가 맨 뒤예요. $A$를 보증하는 지갑은 하나도 없어요. 순간 이동을 빼면, 그 점수는 모두 매달린 점 규칙에서 나와요. 이 규칙은 $D$의 몫을 네 점에 고르게 나눠 줘요($d\,r_D/4=0.1000$). 이 예는 한 점이, 자기도 보증을 받는 점들로부터 보증을 \emph{받아서} 점수를 얻는다는 것을 보여 줘요. 이것이 바로 페이지랭크의 인용 논리예요\cend{15,23}.

같은 지갑들이 허락 대신 송금으로 이어져 있다고 해 봐요. 송금은 보통 다른 쌍들을 잇거나, 같은 쌍이라도 쌍마다 합계가 달라요. 그러면 다른 이동 행렬 $P$가 나오고, 똑같은 점들에 대해서도 순위가 달라질 수 있어요. 가설 H1은 바로 이 상황을 작게 줄여 놓은 것이에요. 엔도스랭크와 AWP가 꼭 일치할 필요는 없어요.

\dissertationSubheading[sec:defi-risk-taxonomy]{디파이 위험의 분류와 평판이 들어갈 자리}

Werner 등\cend{8}은 디파이 위험을 서비스의 층에 따라 나누었고, Zhou 등\cend{27}은 실제로 일어난 공격의 종류에 따라 묶었어요. 계약 코드의 버그(프로그램 오류)와 잘못 설정한 업그레이드 키(계약을 고칠 수 있는 열쇠)는 스마트 계약 위험이에요\cend{50}. 낡은 가격이나 누군가가 마음대로 다루는 가격 정보는 오라클 위험이에요\cend{39}. 시장 위험\cend{51}에는 값이 크게 출렁이는 것과 청산이 줄줄이 이어지는 것이 들어가고, 스테이블코인(값을 일정하게 맞춘 코인) 시장과 자금 조달 시장의 압박도 여기에 들어가요\cend{52}. 2020년 3월의 시장 폭락은 이 두 층이 서로 어떻게 얽힐 수 있는지 보여 줬어요\cend{9}. 모형으로 나타낼 수 있는 토큰 가중 투표(토큰이 많을수록 표가 많은 투표)는 거버넌스(운영 규칙을 정하는 일) 위험에 들어가요\cend{53}. 우리는 더 좁은 층인 \emph{거래 상대} 위험, 곧 \emph{행동} 위험을 연구해요. 이것은 한 주소의 지난 행동을 보고 그 주소가 안전한 거래 상대인지를 묻는 거예요. 레버리지(빌린 돈을 보태 크게 거래하기)를 책임 있게 다루는 트레이더로서든, 많은 주인이 자기 토큰을 쓰도록 믿고 맡기는 계약으로서든 말이에요.

온체인 평판 점수는 행동 위험 층의 일부만 다루고, 담보나 감사(전문가가 코드를 점검하는 일)나 오라클 설계를 대신하지 못해요. 평판 점수의 목적은 한 시스템의 관계 구조(누가 누구에게 허락하는지, 누가 누구에게 돈을 내는지)를, 적은 수고로 확인하고 다시 계산할 수 있는 순서 값으로 압축하는 것이에요\cend{2}. 행동 위험 안에서 엔도스랭크는 허락을 보고, AWP는 돈의 흐름을 봐요. 둘 다 취약점을 노린 공격이나 거버넌스 공격을 예측하려고 하지 않아요.

끝으로, 가짜 활동은 행동 위험을 재기 어렵게 만들어요. 자전 거래는 탈중앙 거래소에서 거래량을 부풀리고\cend{28}, 규제를 받지 않는 중앙 거래소에서는 더 큰 규모로 부풀려요\cend{29}. 앞지르기 거래(front-running, 남의 거래를 보고 먼저 끼어드는 거래)와 그와 비슷한 추출 가능한 가치(거래 순서를 이용해 빼내는 이익)는, 관계가 아니라 순서 전략을 따르는 송금을 만들어요\cend{54}. Qin 등\cend{30}은 이더리움에서 그런 가치가 얼마나 빼내어지는지 어림했어요. 유니스왑(Uniswap)의 사기 토큰\cend{31}과 거래소 사이의 차익 거래(값 차이로 이익을 내는 거래)\cend{55}도 마찬가지로, 오래가는 관계를 만들지 않으면서 거래만 늘려요. 송금 그래프는 가짜 거래량에 특히 약하고, 허락 그래프는 허락 파밍(가짜 허락을 잔뜩 만들어 내기)에 약해요. 시간으로 나누기와 시빌 찾아내기는 이 문제를 일부만 풀어요.

\dissertationSubheading[sec:related-measurement]{관련된 측정 전통}

블록체인 밖에서는 전자 상거래\cend{56}, P2P 네트워크\cend{16}, 웹 스팸 연구\cend{17}에서 평판 시스템을 연구해 왔어요. 이 연구들에서 얻은 네 가지 깨달음이 여기에도 들어맞아요. (i) 평판은 관계가 가진 성질이에요. (ii) 신분을 쉽게 만들 수 있으면 시빌 공격이 가능해져요\cend{11}. (iii) 씨앗 두기와 정규화는 부풀리기의 효과를 줄여요. (iv) 안에서 앞뒤가 맞는 것만으로는 모자라기 때문에, 평가에는 바깥 기준이 필요해요\cend{22}. 엔도스랭크는 (i)을 바로 받아들이고, (ii)는 위협으로 다뤄요. (iii)은 받아들이지 않아요. AWP와 견줄 수 있도록 고른 순간 이동을 그대로 두고, 씨앗 집합을 쓰지 않아요. 그 대가가 얼마인지는 \ref{ch:discussion}장에서 재요. (iv)는 일부만 해내요. 스피어먼의 $\rho$\cend{20}와 켄달의 $\tau$\cend{21}로 재는 같은 기간 점검은 앞뒤가 맞는지를 시험해요. 반면 시간 홀드아웃은 기간 밖 예측을 시험해요. 스펜더에게는 같은 구성개념으로, 트레이더에게는 청산으로 시험해요.

개인화 페이지랭크와 주제 민감 페이지랭크\cend{57,58}는 순간 이동 벡터를 골라 둔 점 무리 쪽으로 기울여요. 이것은 TrustRank가 작은 씨앗 집합을 전체 순위로 바꾸는 데 쓰는 방법이에요\cend{17}. 보내는 활동에 비례하는 AWP의 재시작(다시 출발)도 또 다른 예예요\cend{3}. 다중 페이지랭크(multiplex PageRank)에서는 한 층에서의 점의 중심성이 다른 층의 걷기에 영향을 줘요\cend{62}. 그리고 De Domenico 등의 다층 틀에서는, 걷는 사람 한 명이 서로 이어진 네트워크의 층 안과 층 사이를 오가요\cend{63}. 두 방법 모두 층을 한 점 집합 위의 서로 다른 관계로 다뤄요. 이것은 같은 주소들 사이의 허락과 송금에 딱 맞아요. C-PR은 이 생각의 가장 단순한 꼴을 써요. 곧 점마다 두 이동 행렬을 비율 하나로 섞어 더한 것이에요. S-PR(보증 씨앗 페이지랭크)은 개인화 순간 이동의 생각을 빌려 와서, 허락 점수를 순간 이동 벡터로 써요(\ref{sub:coupled-graph}절).

\dissertationSubheading[sec:account-model]{계정, 가스, 그리고 로그를 재는 바탕으로 삼는 까닭}

이더리움에는 두 종류의 계정이 있어요. 외부 소유 계정(EOA)과 계약 계정이에요. EOA는 개인 키(주인만 아는 비밀 열쇠)가 다루고, 계약 계정은 그 안에 저장된 코드가 다스려요\cend{59}. 상태를 바꾸는 거래는 모두 가스를 써요. 로그는 계약의 \texttt{LOG} 명령에서 나와요. 이 명령은 계약 코드가 이벤트를 내보낼 때(솔리디티(Solidity)의 \texttt{emit} 문장) 실행돼요. 저장 공간에 값을 쓰는 것만으로는 로그가 생기지 않아요. 남이 확인할 수 있는 평판 시스템을 만들려면 오프체인 데이터베이스보다 로그를 골라야 해요. 로그는 블록체인 역사의 영원한 한 부분이라서, 어떤 풀 노드(모든 기록을 가진 컴퓨터)나 색인 제공자(기록을 찾기 쉽게 정리해 주는 서비스)라도 따로 다시 계산할 수 있기 때문이에요\cend{60}.

가스 값은 그래프의 모양에도 영향을 줘요. L2 네트워크에서 가스가 싸지면 $G_{\text{transfer}}$가 $G_{\text{approve}}$보다 더 촘촘해져요. 송금은 늘 되풀이되지만, 허락은 나중의 송금들이 지나갈 수 있도록 한 번만 해 주기 때문이에요. 그 결과로 생긴 그래프 크기는 \ref{ch:results}장에서 보고해요.

\dissertationSubheading[sec:approval-security]{보안과 관련된 사건으로서의 허락}

무제한 허락(\texttt{type(uint256).max})을 하면 주인은 다시 허락하지 않아도 되지만, 그 대신 안전을 내줘야 해요. 허락을 받은 계정이 해킹당하면, 주인이 허락을 다시 정하기 전까지 주인의 토큰을 하나도 빠짐없이 도둑맞을 수 있어요. 그러므로 허락은 기록이 아니라 능력(무언가를 할 수 있는 힘)이에요. 허락을 가진 쪽은 그 허락을 다 쓰거나 거둘 때까지 주인의 토큰을 옮길 수 있어요\cend{18}. 엔도스랭크는 토큰 허락 하나를, 한 스펜더가 그 토큰을 옮길 능력을 보증하는 것으로 풀이해요. 그래서 허락은 송금과 같은 뜻이 아니에요. 둘 다 토큰이 움직이는 일과 관련이 있는데도 그래요. 어떤 계정은 송금을 많이 받으면서도, 남을 대신해 돈을 옮겨도 된다는 허락은 한 번도 받지 못할 수 있어요.

자료에서 무제한 허락은 $2^{256}-1\approx1.16\times10^{77}$ 기본 단위(소수점 조정을 하지 않은 가장 작은 단위)의 금액으로 나타나고, 흔해요(그 개수는 \ref{sub:coupled-graph}절에 있어요). 엔도스랭크의 금액 무게는 무제한 허락도 다른 허락과 똑같이 세요. 하지만 C-PR의 금액을 그대로 쓴 허락 층에서는, 무제한 허락이 그 주인의 줄(행)을 거의 모두 차지해요.

\dissertationSubheading[sec:awp-in-depth]{AWP 자세히 보기: 송금 그래프 기준선이 재는 것}

Do, Do, Nguyen\cend{3}이 제안한 적응형 가중 페이지랭크(AWP)는 거래 그래프의 주소들에 순위를 매겨요. 거래 하나하나는 자기 화살표에, 시간의 로지스틱 함수에 금액의 상한 있는 함수를 곱한 값을 더해요. 그래서 어떤 거래도 화살표 무게에 1보다 많이 더하지 못해요. 걷기는 주소마다 그 주소의 활발함(활동 정도)에 비례해 다시 출발해요. 활발함은 그 주소가 돈을 보낸 주소들마다 가장 최근 거래의 시간 무게를 구해 모두 더한 값이에요. 그다음 순위가 얼마나 정확한지를 들어오는 화살표 수(입차수)와 받은 양의 합에 비추어 검증해요. 이 두 값은 지갑이 끄는 경제적 관심을 대신 재는 대리 지표로 써요. 감쇠의 꼴(식~\eqref{eq:logistic-decay})과, \ref{ch:methodology}장과 \ref{ch:results}장의 입차수/받은 양의 합 대리 지표는 AWP 논문에서 가져왔어요. 그 논문은 매개변수(미리 정해야 하는 값)의 값을 밝히지 않았어요. 여기서 쓴 값은 \ref{sub:awp-graph}절에 있어요. 이 방법의 장점은 모두 온체인이고 순위에 바탕을 둔다는 것이에요. 그리고 어떤 종류의 오프체인 신분도 필요 없어요.

엔도스랭크는 같은 틀을 그대로 써요. 곧 페이지랭크 아래 AWP의 화살표 무게를 쓰고, 순위 상관 점검을 더하는 틀이에요. 다만 화살표가 뜻하는 바를 바꿔요. \ref{ch:results}장에서 우리는 허락 화살표가 구성개념이 다르고, 허락과 잘 맞고, 더 싼 점수를 내는지 시험해요. 또 $t_1$에 매긴 점수로 $t_2$까지의 새 허락을 예측할 수 있는지도 시험해요.

\dissertationSubheading[sec:method-comparison-table]{그래프 평판 방법들의 비교}

표~\ref{tab:method-comparison-lit}에서는 엔도스랭크를 이 장에서 다룬 다른 방법들과, 화살표의 뜻, 시간 모형, 검증 설계를 기준으로 견주어요.

\begin{table}[htbp]
\centering
\caption{골라 뽑은 그래프 평판 방법들의 비교(개념 수준).}
\label{tab:method-comparison-lit}
\fitwidth{%
\begin{tabular}{p{0.18\linewidth}p{0.24\linewidth}p{0.24\linewidth}p{0.28\linewidth}}
\toprule
방법 & 화살표의 뜻 & 시간 모형 & 흔한 검증 방법 \\
\midrule
PageRank (웹) \cend{15} & 하이퍼링크 & 없음 / 정적 스냅숏(멈춘 한 장면) & 검색 품질 \\
EigenTrust \cend{16} & 이웃끼리 매긴 믿음 점수 & 되풀이해서 모으기 & P2P 파일이 진짜인지 \\
TrustRank \cend{17} & 좋은 씨앗에서 나온 링크 & 정적 / 씨앗에서 출발 & 스팸의 순위 낮추기 \\
이더리움의 PageRank\slash HodgeRank \cend{4} & 토큰 송금 & 정적 기간(고정된 한 기간) & 이 논문에서 다시 해 보지 않음 \\
AWP \cend{3} & 거래, 상한 있는 금액 무게 & 화살표와 재시작의 로지스틱 감쇠 & 입차수 / 받은 양의 합 \\
RiskProp \cend{12} & 위험을 실은 연결 & 위험 퍼뜨리기 & 계정 위험 라벨 \\
다중 / 다층 PageRank \cend{62,63} & 한 점 집합 위의 여러 관계 & 정적 층 & 상위 점들의 다재다능함(여러 층에서 두루 중요한 정도) \\
Human Passport \cend{64} & 화살표 없음; 주소 주인에 대한, 무게를 단 온체인·오프체인 자격 증명 & 지금 가진 자격 증명 묶음 & 통과 기준값; 이미 알려진 사람 주소와 시빌 주소 \\
EndorseRank (이 연구) & 마지막 허락, AWP의 무게 & 마지막 허락마다의 로지스틱 감쇠 & 허락 점검; $t_2$까지의 새 허락 \\
C-PR / S-PR (이 연구) & 송금과 섞거나 송금의 씨앗이 되는 허락; 원래 금액 & 허락 스냅숏과 송금의 로지스틱 감쇠 & 홀드아웃 라벨; 미리 정한 규칙 \\
\bottomrule
\end{tabular}
}
\end{table}

표~\ref{tab:method-comparison-lit}에서 보듯이, 엔도스랭크의 기여는 순위 매기는 알고리즘이 아니라, 마지막 허락을 보증 화살표로 고른 데에 있어요. 연산자 쪽에서 우리의 기여는 두 관계를 하나의 걷기로 결합한 것이에요. 여기에 씨앗을 쓴 빼 보기 실험(한 부분을 바꿔서 그 몫을 재는 실험)을 더해, 이 결합이 얼마나 보탬이 되는지 재요. 결합은 기간 밖에서, 두 구성개념 모두에 대해, 미리 정한 규칙 아래에서 시험해요.

\dissertationSubheading[sec:literature-summary-table]{살펴본 문헌의 요약}

표~\ref{tab:literature-summary}에는 이 연구의 바탕이 되는 자료들을 한곳에 모았어요. 각 줄에는 자료의 연구 방법, 이 논문이 가져다 쓰는 장점, 이 논문이 해결하려는 한계, 그리고 그 자료가 이 연구와 어떻게 이어지는지를 적었어요. \ref{sec:research-gap}절에서 밝힌 연구의 빈틈은 이 표에 모은 증거에 바탕을 둬요.

{\footnotesize
\setlength{\tabcolsep}{4pt}
\begin{longtable}{p{0.15\linewidth}p{0.20\linewidth}p{0.19\linewidth}p{0.19\linewidth}p{0.18\linewidth}}
\caption{살펴본 문헌의 요약: 연구 방법, 장점, 한계, 이 연구와의 관계.}
\label{tab:literature-summary}\\
\toprule
자료 (연도) & 연구 방법 & 장점 & 한계 & 이 연구와의 관계 \\
\midrule
\endfirsthead
\multicolumn{5}{l}{\small 표~\thetable\ (이어짐)}\\
\toprule
자료 (연도) & 연구 방법 & 장점 & 한계 & 이 연구와의 관계 \\
\midrule
\endhead
\midrule
\multicolumn{5}{r}{\small 다음 쪽에 이어짐}\\
\endfoot
\bottomrule
\endlastfoot
Page et al.\ (1998) \cend{15}; Langville과 Meyer (2006) \cend{23} & 하이퍼링크 그래프에 고유벡터로 순위 매기기, 거듭제곱 반복으로 풀기 & 해석할 수 있는 퍼뜨림; 반복 한 번에 $O(|E|)$; 감쇠 덕분에 페론 벡터가 하나로 정해짐 & 화살표는 하이퍼링크만 뜻할 수 있음; 시간 개념이 없음 & 같은 풀이 프로그램과 복잡도 논리를, 비교하는 두 방법에 바꾸지 않고 그대로 씀 \\
Kamvar et al.\ (2003) \cend{16}; Gy\"{o}ngyi et al.\ (2004) \cend{17} & P2P 네트워크에서 믿음 모으기; 스팸을 끌어내리려고 씨앗에서 퍼뜨리기 & 정규화와 씨앗 집합이 부풀리기를 막음; 공격자를 분명하게 모형에 넣음 & 직접 매긴 점수나 골라 둔 씨앗에 기댐; 온체인에 쓰인 적 없음 & 시빌 비용 논의와 고른 순간 이동 기준선에 도움을 줌 \\
Do와 Do (2023) \cend{4} & 사회적 신용을 재려고 이더리움 거래 그래프에 PageRank와 HodgeRank를 씀 & 송금으로 이어진 정도를 신용과 관련된 것으로 봄; 자료가 모두 공개됨 & 송금은 돈 내기와 믿음을 뒤섞음; 허락의 뜻은 없음 & 앞선 송금 그래프 연구, 이 논문에서는 AWP와 따로 다룸 \\
Do, Do, Nguyen (2023), AWP \cend{3} & 상한 있는 금액 무게, 로지스틱 최근성(최근 것일수록 무겁게 세기), 보낸 활동으로 무게를 단 재시작을 쓰는 거래 PageRank; 검증은 입차수/받은 양의 합과의 순위 상관 & 순위로 검증하므로 채무 불이행 라벨이 필요 없음; 최근성을 분명하게 모형에 넣음 & 대리 지표를 그래프를 만든 바로 그 송금 자료로 계산함(저절로 생기는 일치도); 촘촘한 그래프; 매개변수 값을 밝히지 않음 & 기준선, 발표된 그대로 씀(\ref{sub:awp-graph}절); 감쇠 매개변수와 $b$는 여기서 정하고, EndorseRank도 같은 화살표 무게를 씀 \\
Lin et al.\ (2025), RiskProp \cend{12} & 계정 위험을 연결망으로 퍼뜨리기, 익명 벗기기 점수와 짝지음 & 구조가 활동 개수보다 더 많은 정보를 보탠다는 것을 보임; 라벨로 평가함 & 위험 라벨이 필요함; 믿음이나 허락이 아니라 사기를 겨냥함 & 화살표의 뜻을 더 섬세하게 나눠야 한다는 근거; 청산으로 무게를 단 단순한 송금 걷기는 보관해 둠(부록~\ref{ch:appendix-archive}) \\
Wu et al.\ (2022) \cend{24} & 이더리움 거래 그래프의 연결망 임베딩으로 피싱 찾기 & 그래프 특징이 계정 하나하나의 특징보다 나음; 라벨이 붙은 큰 자료 & 지도 학습(정답을 보며 배우기); 라벨은 신용이 아니라 피싱을 표시함 & 그래프 속 위치가 정보를 담는다는 전제를 뒷받침함; 경쟁하는 방법은 아님 \\
Hassija et al.\ (2020) \cend{13}; Kotb (2023) \cend{14} & 전망 이론과 기계 학습으로 블록체인 자료의 신용 점수 매기기 & 대출 결정을 바로 겨냥함 & 라벨이나 오프체인 특징이 필요함; 일부만 해석할 수 있음 & 대비되는 부류; EndorseRank는 확인할 수 있는 그래프 퍼뜨림 안에 머묾 \\
Bellini et al.\ (2020) \cend{25}; Almasoud et al.\ (2020) \cend{26} & 블록체인 기반 믿음·평판 시스템 조사 & 설계 공간을 지도처럼 정리하고, 직접 매긴 점수와 거래량 신호가 대부분임을 밝힘 & 허락 화살표는 평가하지 않음; L2에서 잰 것이 없음 & 조사한 설계 가운데 마지막 허락 화살표를 쓰는 것이 하나도 없음을 보여 줌 \\
Victor (2020) \cend{40}; Victor와 Weintraud (2021) \cend{28}; Cong et al.\ (2023) \cend{29} & 주소를 묶는 어림 규칙; 탈중앙 거래소(DEX)와 거래소에서 자전 거래 찾기 & 온체인 활동 가운데 꾸며 낸 것이나 여러 주소로 한 것의 비율을 잼 & 평판 점수는 내놓지 않음 & 송금 신호가 왜 쉽게 부풀려지는지 보여 줌, 이것이 허락 구성개념을 쓸 근거가 됨 \\
Cronbach와 Meehl (1955) \cend{22}; Messick (1995) \cend{33} & 심리 측정학의 구성 타당도(점수가 재려는 것을 정말 재는지) & 점수가 재도록 설계된 것과 바깥 기준을 구별함 & 블록체인만을 위한 것은 아님 & 같은 기간 점검을 구성개념 점검으로, 시간 홀드아웃을 기간 밖 증거로 다루는 바탕 \\
Human Passport (연도 없음) \cend{64}; Siddarth et al.\ (2020) \cend{65}; Bordeianu와 Popescu (2026) \cend{66} & 주소 주인에 대한, 무게를 단 온체인·오프체인 자격 증명을 통과 기준값과 견줌; 사람 증명 방어에 대한 검토 & 큰 규모로 실제 쓰임; 시빌 공격을 막으려고 오프체인 증거를 씀 & 순서를 매기는 점수가 아니라 문; 주소 사이의 관계는 아무 역할도 없음; 자격 증명 파밍이 남아 있음 & 실제로 믿음을 어떻게 쌓는지, 그리고 EndorseRank가 주지 못하는 것을 보여 줌: 사람임을 보장하지 않음 \\
Efron (1979) \cend{34}; DiCiccio와 Efron (1996) \cend{35} & 부트스트랩 다시 뽑기와 부트스트랩 신뢰구간 & 분포를 가정하지 않고 순위 통계의 불확실성을 구함; 짝지은 대비 & 단위들이 서로 바꿔도 되는 것이어야 함(교환 가능성) & \ref{ch:results}장에서 보고하는 95\% 구간과 $\Delta\tau$ 대비의 출처 \\
\end{longtable}
}

\dissertationSubheading[sec:research-gap]{연구의 빈틈}

표~\ref{tab:literature-summary}와 \ref{sec:related-measurement}절에서 다섯 가지 빈틈이 드러나요. 첫째, 여기서 다룬 온체인 평판 방법은 모두(PageRank/HodgeRank \cend{4}, AWP \cend{3}, RiskProp \cend{12}, 조사 연구에 나온 설계들 \cend{25}) 송금이나 직접 매긴 점수로 그래프를 만들어요. 그리고 디파이에서 쓰이는 증명 도구는 주소 사이의 관계 그래프 없이 주소를 가진 사람에게 점수를 매겨요\cend{64}. 어느 것도 ERC-20 허락을 쓰지 않아요. 이 허락은 모든 ERC-20 토큰에 담긴, 쓸 힘을 일부러 맡긴 일의 온체인 기록이에요. 마지막 허락으로 만든 보증 그래프가 다르고 뜻있는 순서를 내는지는 아직 알려지지 않았어요(O1과 RQ1). 둘째 빈틈은 검증에 있어요. 송금 그래프 기준선은 그래프를 만들 때 쓴 바로 그 송금 자료로 계산한 입차수와 받은 양의 합에 점수를 견주어 점검해요. 그래서 어느 정도의 일치는 피할 수 없어요. 같은 주소 묶음에서, 허락 구성개념과 그 짝인 송금 구성개념에 이 절차를 신뢰구간 및 짝지은 차이 검정과 함께 적용한 보고는 없어요(O2와 RQ2). 셋째, 블록체인 페이지랭크의 실행 시간은 정해진 절차에 따른 화살표 수, 메모리 사용량, 반복 횟수 없이 발표되는 일이 많아요. 그래서 효율이 좋다는 주장을, 구현 방식이 아니라 희소성 덕분이라고 말하기 어려워요(O3과 RQ3). 넷째 빈틈은 방법에 관한 것이고, 시간과 관련이 있어요. 같은 기간의 그래프 통계가 검증처럼 제시되는 일이 많지만, 그런 점수가 앞으로의 행동을 예측한다는 것은 보일 수 없어요. 그 대신 여기서 쓰는 기간 밖 점검은 시간 홀드아웃이에요. 구성개념마다 앞으로 일어날 일을 세는 라벨을 두고, 동결 때의 원시 차수를 기준선으로 써요(O4와 RQ4). 끝으로, 다층 순위 매기기는 사회 연결망, 교통망, 인용 연결망을 위해 발전해 왔지만\cend{62,63}, 한 주소 묶음 위의 두 온체인 관계에는 쓰인 적이 없어요. 허락 층과 돈 내기 층을 합치면 돈 내기 층 하나만 쓸 때보다 나은지, 또는 두 층을 모두 이기는지를, 결과를 보기 전에 정한 규칙 아래에서 보고한 연구는 없어요(O5와 RQ5).

\dissertationSubheading[sec:open-problems]{이 연구를 이끈 풀리지 않은 문제들}

앞선 문헌은 온체인 평판에 대한 몇 가지 문제를 풀지 못한 채 남겨 두었어요. 이 문제들이 이번 연구를 하게 된 까닭이에요:

\begin{enumerate}
\item 화살표의 뜻: 신분이 가명일 때, 어떤 온체인 관계가 믿음이나 허락을 가장 잘 담을까요\cend{2}?
\item 채무 불이행 없이 하는 검증: 대출을 못 갚았다는 라벨이 드물 때, 무엇을 바탕으로 방법들을 비교할 수 있을까요\cend{4}?
\item 계산 비용: 그래프가 커져도, 평판 점수를 사람이 기다리지 않아도 될 만큼 빨리 새로 고칠 수 있을까요\cend{23}?
\item 시간을 나눈 점검: 평가가 같은 기간의 그래프 통계를 기간 밖 증거로 잘못 아는 일을, 무엇으로 막을 수 있을까요\cend{22}?
\item 공격에 버티는 힘: 신분을 싸게 만들 수 있을 때, 평판을 꾸며 내는 데 비용이 얼마나 들까요\cend{11}?
\end{enumerate}

엔도스랭크는 허락 화살표, 같은 기간 점검, 실행 시간 벤치마크(속도 재기), 시간 홀드아웃으로 (1)--(4)를 다뤄요. 문제 (5)는 \ref{ch:discussion}장에서 모형 수준으로 분석해요. 실제 자료로 공격을 온전히 평가하는 일은 앞으로의 연구로 남겨 둬요.

\ref{ch:methodology}장에서는 이 빈틈 이야기를 실제로 할 수 있는 설계로 바꿔요. 그 설계에는 자료 기간과 코호트, 두 그래프를 만들고 하나의 공통 페이지랭크 구현으로 푸는 방법, 대리 지표 가족, 시간 홀드아웃, 그리고 \ref{ch:results}장 결과의 바탕이 되는 시간 재기 절차와 부트스트랩 절차가 들어 있어요.

